\chapter{110}



«Danke, Maxime, und lass mir deine Mutter grüssen.»

Manuela winkte dem Postboten nach. Und noch während das Elektrofahrzeug
wendete, suchte sie nach dem Absender auf dem Paket. Ein Fotodienst? Was
Joh da wohl bestellt hatte? Beim genaueren Hinsehen war sie verwirrt. Das
Paket war nicht an Joh adressiert. \emph{Ava Brand}, las sie, und kaum
hatte sie es auf den Tisch gelegt, stürmte Ava in die Küche.

«Ich habe Maxime gesehen, es ist da,» jauchzte sie, setzte sich an den
Tisch und begann das Paket hektisch zu öffnen. Manuela schaute ihr
vergnügt zu.

«Dann ist das also für dich?»

Die perforierte Kartonlasche zwischen Avas Fingernägeln liess sich
schwer aufreissen. Sie verzerrte das Gesicht.

«Was für eine dämliche Verpackung.»

Sie setzte erneut an, presste die Lippen aufeinander, zerrte und fluchte
und schaffte es schliesslich. Der Karton liess sich auseinander klappen.

«Ja und nein, das ist für uns alle.»

Ava wirkte, als würde sie nächstens einen Freudentanz aufführen.
Manuela setzte sich ihr gegenüber an den Tisch. Joh hatte es neulich treffend
ausgedrückt. Die junge Frau hatte sich eine rührende Kindlichkeit bewahrt.
Manuela beugte sich vor.

«Oh, ich bin sowas von gespannt. Was ist denn da drin?»

«Hm, warte,» Ava legte rasch ihre Arme über den Inhalt, «ich überlege,
wie ich es machen soll.»

Sie blickte zur Decke, rollte die Augen in alle Richtungen und setzte
sich mit einem Mal kerzengerade auf.

«Gut, ich, ich zeige dir, was es ist, aber nur von aussen. Ansehen
werden wir es uns, wenn alle da sind.»

Sie schmunzelte und hob einen in Seidenpapier eingeschlagenen, schweren Gegenstand aus dem Karton.
Manuela rückte ein Stück näher und setzte sich ihre Lesebrille auf.
Ava löste das Papier und ein Buch kam zum Vorschein. Augenblicklich deckte sie den untern Teil der Frontseite mit ihren Händen ab. 
 
«Das ist ein, ein Fotobuch,» Manuelas Staunen liess Ava geheimnisvoll
nicken.

«Ja, genau.»

«\emph{Schafscheid}, 21./22. September, Unsere Schafe kehren heim», las
Manuela, «ein Fotobuch zur \emph{Schafscheid}? Hast du das gemacht?»

«Mhm,» Ava grinste, «es sollte eine Überraschung werden. Ich glaube,
keiner hat es bemerkt. Es wurde so viel geknipst, dass ich nicht
aufgefallen bin.»

Sie sah stolz aus. Jetzt schob sie ihre Hände ganz langsam nach unten
und das Titelbild wurde sichtbar.

«Nein, ich glaub’s ja nicht, Joh und ich. Und \emph{Perla},» Manuela war ausser sich, «und ihr Lämmchen, du bist verrückt, Ava. Das ist  unglaublich.»

Manuela stand auf, umrundete den Tisch und umarmte Ava von hinten. Die
fasste nach ihren Händen und blickte schräg zu ihr hoch.

«Schön, dass du dich freust. Ich möchte mich mit dem Buch bei dir und
Joh bedanken. Die \emph{Schafscheid} war ein einziges Abenteuer für
mich, auch weil Dad mit dabei war. Und die Geburten der Lämmer
mitzuerleben, Manuela, war so unendlich beeindruckend, und über die Bilder der Schafschur freue ich mich am meisten. Ich
glaube, die sind mir am besten gelungen.»

«Ich bin ja so was von gespannt, Ava.»

Ava stand auf, drehte sich um und fasste Manuela bei den Händen.

«Und was ich dir schon länger sagen möchte, Manuela,» der Ausdruck in
ihrem Gesicht veränderte sich unverkennbar, «es ist so wertvoll, dass ich hier oben
mit euch wieder ins Leben zurückfinde. Ich schätze die Lebendigkeit,
sogar den Trubel, bei euch. Und immer wieder bin ich fasziniert, wie ihr
Problemen begegnet. Manchmal bin ich richtig glücklich und zwischendurch
vergesse ich sogar Daniel.»

Manuela lächelte und zog sie an sich. Ava kämpfte. 
Dann schien sie sich wieder gesammelt zu haben.

«Was meinst du, wann sind alle da?»

Sie wischte sich einzelne Tränen aus den Augen.
Manuela legte den Kopf schief.

«Hm, Buchvernissage nach dem Nachtessen?»
